\documentclass[border=10pt]{standalone} 
\usepackage[utf8]{inputenc}
\usepackage{nacal}
\usepackage{pgfplots}
\usepackage{tikz-3dplot}
\usepackage{tikz} 
\usetikzlibrary{matrix,decorations.pathreplacing}
\usetikzlibrary{calc,angles,positioning,intersections,quotes,decorations.markings}
\usepackage{tkz-euclide}

\pgfplotsset{width=10cm,height=6cm}

\begin{document}
  \def\Rojo{red!50!black}
  \def\RojoOscuro{red!20!black}
  \def\RojoClaro{red!90!black}
  \def\Verde{green!50!black}
  \def\VerdeOscuro{green!30!black}
  \def\VerdeClaro{green!50!gray}
  \def\Azul{blue!50!black}
  \def\AzulOscuro{blue!35!black}
  \def\AzulClaro{blue!60!gray}

  \newcommand{\Prho}{.9}%
  \newcommand{\Ptheta}{55}%
  \newcommand{\Pphi}{60}%  
  \tdplotsetmaincoords{90}{90}
  \begin{tikzpicture}
    [scale=6,
     tdplot_main_coords,
     axis/.style={-,gray},
     vector/.style={-{Stealth[length=3mm, width=2mm]},\Verde,very thick},
     guide/.style={dashed,thin, gray},
     vector guide/.style={-{Stealth[length=2mm, width=1mm]},dashed,\AzulClaro,thick}]
     
    %standard tikz coordinate definition using x, y, z coords
    \coordinate (O) at (0,0,0);
    %tikz-3dplot coordinate definition using r, theta, phi coords
    \tdplotsetcoord{P}{\Prho}{\Ptheta}{\Pphi}
    %draw axes
    %\draw[axis] (0,0,0) -- (0.65,0,0); %  node[anchor=north east]{$x$};
    \draw[axis] (0,0,0) -- (0,.9,0); % x1 node[anchor=north west]{$y$};
    \draw[axis] (0,0,0) -- (0,0,.6); % x2 node[anchor=south]{$z$};
    %draw a vector from O to P
    \draw[vector] (O) -- (P) node[right] {$\Vect{x}=(x_1,x_2,)$};
    %draw guide lines to components
    \draw[vector guide] (O) -- (Pxy) node [below] {$x_1$};
    \draw[vector guide] (O) -- (Pz)  node [left]  {$x_2$};
    % Compute x,y,z
    \pgfmathsetmacro{\PxCoord}{\Prho * sin(\Pphi) * cos(\Ptheta)}%
    \pgfmathsetmacro{\PyCoord}{\Prho * sin(\Pphi) * sin(\Ptheta)}%
    %\pgfmathsetmacro{\PzCoord}{\Prho * cos(\Pphi)}%
    \draw[guide] (P) -- (Pz); %{\PxCoord};
    \draw[guide] (P) -- (Pxy);
    %\node[below,black]  at (Pxy) {$(x_1,x_2,)$};                  
  \end{tikzpicture}
\end{document}
